
\section{\trTransversality}
\label{sec:transversality}

我们现在就来证明\trTransversality{}性质（命题~8.1.4），证明中要用到 $(d\scr{F})_{u}$ 的\trFredholm{}性质，也就是说，这里先把定理~8.1.5 的第一部分当作已知的而暂置不论。\trAlmostComplex{}结构 $J$ 是取定的，\trNondegenerate{}的\trHamiltonian{} $H_{0}$ 也是如此。我们取定一个 $h_{0}\in\cC^{\infty}_{\varepsilon}(H_{0})$，以及\trFloerEquation{}（对 $J$ 与 $H=H_{0}+h_{0}$）的一个解 $u$：
\[
\scr{F}(u)=\frac{\partial u}{\partial s}+J\frac{\partial u}{\partial t}+\operatorname{grad}_{u}(H)=0\text{。}
\]
$\scr{F}$ 的\trDifferential，正如我们刚刚解释过的，是在沿 $u$ 取定的一个酉\trTrivialization{}中计算的；我们把它写作
\[
(d\scr{F})_{u}(Y)=\frac{\partial Y}{\partial s}+J_{0}\frac{\partial Y}{\partial t}+S\cdot Y\text{。}
\]

\subsection{$\Gamma$ 的定义}
\label{subsec:definition-gamma}

我们取定两个不同的\trCriticalPoint{} $x$ 与 $y\in\operatorname{Crit}\scr{A}_{H_{0}}$。为了证明 $\scr{Z}(x,y,J)$ 是一个\trManifold，我们把它描述成某个\trFiberBundle{}的某个\trSection{}的零点所成之集。在空间 $\scr{P}(x,y)\times\cC^{\infty}_{\varepsilon}(W,H_{0})$ 上，\footnote{译注：原书此处印作 $\cC^{\infty}_{\varepsilon}(W,H_{0})$，多出一参数 $W$——本节前后各处（紧随其下的两次即然）皆作 $\cC^{\infty}_{\varepsilon}(H_{0})$，疑系衍余，此处照原书排印，存疑备案。——译者注}我们定义\trVectorBundle
\[
\scr{E}=\bigl\{(u,h,Y)\mid Y\in L^{p}(u^{\star}TW)\bigr\}\text{。}
\]
我们可以把 $\scr{E}$ 看成一个\trFiberBundle
\[
\scr{E}\longrightarrow\scr{P}(x,y)\times\cC^{\infty}_{\varepsilon}(H_{0})
\]
（通过 $(u,h,Y)\mapsto(u,h)$）。映射
\[
\begin{aligned}
\scr{P}(x,y)\times\cC^{\infty}_{\varepsilon}(H_{0})&\ \xrightarrow{\ \sigma\ }\ \scr{E}\\
(u,h)&\longmapsto\frac{\partial u}{\partial s}+J(u)\frac{\partial u}{\partial t}+\operatorname{grad}_{u}(H_{0}+h)
\end{aligned}
\]
是这个\trFiberBundle{}的一个\trSection。\trSection{} $\sigma$ 确实取值于 $\scr{E}$ 之中——也就是说，
\[
\frac{\partial u}{\partial s}+J\frac{\partial u}{\partial t}+\operatorname{grad}_{u}(H)
\]
确实属于 $L^{p}(\RR\times S^{1};u^{\star}TW)$——这一事实，正如我们已在 §8.2.d 中指出过的，将在第 13 章中得到证明（更准确地说，是在引理~13.3.1 与注记~13.3.2 中）。

$\sigma$ 的零点恰好就是 $\scr{Z}(x,y,J)$ 的那些点（更确切地说，是其 $W^{1,p}$ 版本的点）；注意，数据 $(H_{0},J)$ 现在是 $\cC^{\infty}$ 的，而 $(u,Y)$ 则是\trWeakSolution。要得出结论，只需证明 $\sigma$ 与\trZeroSection\trTransverse，也就是说，当 $\sigma(u,h)=0$ 时，\trLinearization{}后的映射 $(d\sigma)_{(u,h)}$ 与到\trFiber{}上的投影之复合是\trSurjective{}的。

注意，$\sigma$ 关于 $h$ 是线性的，而它的\trDifferential{}为
\[
(d\sigma)_{(u,h)}(Y,\eta)=(d\scr{F})_{u}(Y)+\operatorname{grad}_{u}(\eta)\text{。}
\]
映射
\[
\Gamma\colon W^{1,p}(\RR\times S^{1};\RR^{2n})\times\cC^{\infty}_{\varepsilon}(H_{0})\longrightarrow L^{p}(\RR\times S^{1};\RR^{2n})\text{，}
\]
其\trSurjective{}性正是命题~8.1.4 所断言的，它是 $(d\sigma)_{(u,h)}$ 与从 $T_{\sigma(u,h)}\scr{E}$ 到 $\scr{E}$ 在 $(u,h)$ 处的\trFiber{}上的\trTangentSpace{}之投影的复合。于是，命题~8.1.3 立即从命题~8.1.4 推出。注意，$TW$ 沿 $u$ 的\trTrivialization{} $Z_{i}$ 使我们得以把这个\trFiber{}的\trTangentSpace{}与空间 $L^{p}(\RR\times S^{1};\RR^{2n})$ 等同起来。

\subsection{命题~8.1.4~的证明}
\label{subsec:proof-prop-814}

假设 $\Gamma$ 不是\trSurjective{}的。这时我们就有下面的结果。

\begin{lemma}\label{thm:z-pairing-vanish}
设\footnote{译注：由 $1/p+1/q=1$（本节 $p>2$）即得 $q=p/(p-1)>1$——共轭指数关系自动给出 $q>1$，不碍下文对 $q$ 的使用。——译者注} $q>0$ 满足 $1/p+1/q=1$。那么，存在一个 $\cC^{\infty}$ 类的非零\trVectorField{} $Z\in L^{q}(\RR\times S^{1};\RR^{2n})$，使得对每个 $h\in\cC^{\infty}_{\varepsilon}(H_{0})$ 与每个 $Y\in W^{1,p}(\RR\times S^{1};\RR^{2n})$，都有
\[
\langle Z,(d\scr{F})_{u}(Y)\rangle=0\qquad\text{以及}\qquad\langle Z,\operatorname{grad}_{u}(h)\rangle=0\text{。}
\]
\end{lemma}

在这个陈述中，$\langle\ ,\ \rangle$ 表示乘积（即 $L^{q}$ 与 $L^{p}$ 的配对）
\[
\langle Z,Y\rangle=\int_{\RR\times S^{1}}\langle Z(s,t),Y(s,t)\rangle\,ds\,dt
\]
而且，与上面一样，我们已经借助一个取定的酉\trTrivialization{}把 $W^{1,p}(\RR\times S^{1};u^{\star}TW)$ 与 $W^{1,p}(\RR\times S^{1};\RR^{2n})$ 等同了起来（对 $L^{q}$ 亦然）。

\begin{proof}[\cref{thm:z-pairing-vanish}的证明]
我们先来证明，$\Gamma$ 的\trImage{}是 $L^{p}(\RR\times S^{1};\RR^{2n})$ 的一个闭子空间。算子 $\Gamma$ 不是\trFredholmOperator{}（参见第 14 章的练习~44），但 $(d\scr{F})_{u}$ 是——这将在\cref{subsec:proof-prop-871}中证明。于是，它的\trImage{}是闭的且余维数有限。当然，这个\trImage{}含于 $\Gamma$ 的\trImage{}之中。因此，存在一个有限维的子空间 $G$，它是 $\mathrm{Im}(d\scr{F})_{u}$ 的\trComplement，也就是说，
\[
\mathrm{Im}(d\scr{F})_{u}\oplus G=L^{p}(\RR\times S^{1};\RR^{2n})
\]
并且有连续的投影\footnote{必要时，参见 [17,~Chapter~2]。}
\[
\pi_{1}\colon L^{p}(\RR\times S^{1};\RR^{2n})\longrightarrow\mathrm{Im}(d\scr{F})_{u}\qquad\text{以及}\qquad\pi_{2}\colon L^{p}(\RR\times S^{1};\RR^{2n})\longrightarrow G\text{。}
\]
子空间
\[
G_{1}=G\cap\mathrm{Im}(\Gamma)\subset G
\]
在 $G$ 中维数有限，而 $G$ 本身也是有限维的；因此它有一个\trComplement{} $G_{2}$。于是我们有
\[
\mathrm{Im}(\Gamma)=\mathrm{Im}(d\scr{F})_{u}\oplus G_{1}\text{，}\qquad\text{从而 }\ \mathrm{Im}(\Gamma)\oplus G_{2}=L^{p}(\RR\times S^{1};\RR^{2n})\text{。}
\]
因而存在一个连续投影
\[
\pi_{2}^{\prime}\colon L^{p}(\RR\times S^{1};\RR^{2n})\longrightarrow G_{2}\text{。}
\]
由此可见，$\mathrm{Im}(\Gamma)$ 是闭的：如果由 $\mathrm{Im}(\Gamma)$ 的点组成的一个序列 $(y_{h})$ 在 $L^{p}(\RR\times S^{1};\RR^{2n})$ 中\trConvergence{}于极限 $y_{\star}$，那么 $\pi_{2}^{\prime}(y_{\star})=0$，从而 $y_{\star}\in\mathrm{Im}(\Gamma)$。

于是，$\Gamma$ 的\trImage{}是闭的，而由\trHahnBanach{}定理\footnote{必要时，再次参见 [17,~Chapter~1]。}，存在 $L^{p}(\RR\times S^{1};\RR^{2n})$ 上的一个非零连续线性形式 $\varphi$，它在 $\mathrm{Im}(\Gamma)$ 上为零。\trRiesz{}表示定理\footnote{这一次是 [17,~Chapter~4]。}断言，$\varphi$ 具有形式
\[
\varphi(Y)=\langle Z,Y\rangle
\]
其中 $Z\in L^{q}(\RR\times S^{1};\RR^{2n})$ 是一个\trVectorField。

于是，对每个 $Y\in W^{1,p}(\RR\times S^{1};\RR^{2n})$ 与每个 $h\in\cC^{\infty}_{\varepsilon}(H_{0})$，都有
\[
\langle Z,(d\scr{F})_{u}(Y)\rangle=0\qquad\text{以及}\qquad\langle Z,\operatorname{grad}_{u}(h)\rangle=0\text{。}
\]
我们还须证明这个\trVectorField{} $Z$ 是 $\cC^{\infty}$ 类的。设
\[
L\colon W^{1,p}(\RR\times S^{1};\RR^{2n})\longrightarrow L^{p}(\RR\times S^{1};\RR^{2n})
\]
表示算子 $(d\scr{F})_{u}$。回忆（\cref{thm:df-dbar-plus-s}）
\[
LY=\frac{\partial Y}{\partial s}+J_{0}\frac{\partial Y}{\partial t}+S\cdot Y
\]
（这里 $Y\in u^{\star}TW$；本应使用\trOrthonormal{}\trFrame{} $(Z_{i})$ 在 $W^{1,p}(\RR\times S^{1};\RR^{2n})$ 上定义 $L$，但我们避免这样做，以免加重阅读的负担）。我们如下定义一个映射 $L^{\star}$：
\[
\begin{aligned}
L^{\star}\colon W^{1,q}(\RR\times S^{1};\RR^{2n})&\longrightarrow L^{q}(\RR\times S^{1};\RR^{2n})\\
X&\longmapsto-\frac{\partial X}{\partial s}+J_{0}\frac{\partial X}{\partial t}+{}^{t}S\cdot X\text{。}
\end{aligned}
\]
显然，$L^{\star}$ 是一个\trBounded{}算子。它是 $L$ 的\textbf{\trAdjointOperator}（adjoint operator）：用一个与下面所给的证明相类似的证明，我们可以证明，对 $Y\in W^{1,p}(\RR\times S^{1};\RR^{2n})$ 与 $X\in W^{1,q}(\RR\times S^{1};\RR^{2n})$，有
\[
\int_{\RR\times S^{1}}\langle LY,X\rangle=\int_{\RR\times S^{1}}\langle Y,L^{\star}X\rangle\text{。}
\]
于是，\trVectorField{} $Z$ 理应满足 $L^{\star}Z=0$……但\trApriori{}地看，它取值于 $L^{q}(\RR\times S^{1};\RR^{2n})$ 之中，这比 $L^{\star}$ 的定义空间要大。因此让我们来证明下面的结果。

\begin{lemma}\label{thm:lstar-z-distribution}
按\trDistribution{}的意义，$L^{\star}Z=0$。
\end{lemma}

\begin{proof}
只需对 $Y\in\cC^{\infty}_{0}(\RR\times S^{1};\RR^{2n})$ 与 $X=Z$ 建立上面的等式\footnote{译注：原书本句印作 $Y\in\cC^{0}_{\infty}(\RR\times S^{1};\RR^{2n})$（上标 $0$、下标 $\infty$）；证明明言假设 $\lim_{s\to\pm\infty}Y=0$，所需即具有紧\trSupport{}的\trSmooth{}函数，疑为 $\cC^{\infty}_{0}$ 之上下标互换，此处径改。——译者注}。我们先写
\[
\begin{aligned}
\int_{-\infty}^{+\infty}\int_{0}^{1}\langle LY,Z\rangle\,dt\,ds={}&\int_{0}^{1}\left(\int_{-\infty}^{+\infty}\Big\langle\frac{\partial Y}{\partial s},Z\Big\rangle ds\right)dt\\
&+\int_{-\infty}^{+\infty}\int_{0}^{1}\Big\langle J_{0}\frac{\partial Y}{\partial t},Z\Big\rangle dt\,ds+\int_{-\infty}^{+\infty}\int_{0}^{1}\langle SY,Z\rangle\,dt\,ds\text{。}
\end{aligned}
\]
让我们逐一来考虑这三个积分。首先，由一次简单的分部积分（实际上是\trWeakDerivative{}的性质），
\[
\int_{0}^{1}\left(\int_{-\infty}^{+\infty}\Big\langle\frac{\partial Y}{\partial s},Z\Big\rangle ds\right)dt=-\int_{0}^{1}\left(\int_{-\infty}^{+\infty}\Big\langle Y,\frac{\partial Z}{\partial s}\Big\rangle ds\right)dt
\]
（注意这是成立的，或者朴素地说，那个所谓的“完全积出的”项为零，这是因为我们已经假设了 $\lim_{s\to\pm\infty}Y(s,t)=0$）。对于第二个积分，我们有
\[
\begin{aligned}
\int_{-\infty}^{+\infty}\int_{0}^{1}\Big\langle J_{0}\frac{\partial Y}{\partial t},Z\Big\rangle dt\,ds&=-\int_{-\infty}^{+\infty}\int_{0}^{1}\Big\langle\frac{\partial Y}{\partial t},J_{0}Z\Big\rangle dt\,ds\\
&=\int_{-\infty}^{+\infty}\int_{0}^{1}\Big\langle Y,\frac{\partial J_{0}Z}{\partial t}\Big\rangle dt\,ds=\int_{-\infty}^{+\infty}\int_{0}^{1}\Big\langle Y,J_{0}\frac{\partial Z}{\partial t}\Big\rangle dt\,ds
\end{aligned}
\]
（这一次，“完全积出的”项由关于 $t$ 的\trPeriodic{}性而为零）。接着，我们有
\[
\int_{-\infty}^{+\infty}\int_{0}^{1}\langle SY,Z\rangle\,dt\,ds=\int_{-\infty}^{+\infty}\int_{0}^{1}\langle Y,{}^{t}SZ\rangle\,dt\,ds\text{。}
\]
总括起来，
\[
\begin{aligned}
\int_{-\infty}^{+\infty}\int_{0}^{1}\langle LY,Z\rangle\,dt\,ds&=\int_{-\infty}^{+\infty}\int_{0}^{1}\Big\langle Y,-\frac{\partial Z}{\partial s}+J_{0}\frac{\partial Z}{\partial t}+{}^{t}SZ\Big\rangle ds\,dt\\
&=\langle Y,L^{\star}Z\rangle\qquad\text{（按\trDistribution{}的意义）。}
\end{aligned}
\]
因此，$\langle Z,LY\rangle=0$ 蕴含
\[
\int_{-\infty}^{+\infty}\int_{0}^{1}\langle Y,L^{\star}Z\rangle\,dt\,ds=0\text{，}
\]
换句话说，按\trDistribution{}的意义有 $L^{\star}Z=0$，这正是所要证明的。
\end{proof}

于是，由\trEllipticRegularity{}（定理~12.1.3）推出，$Z$ 是 $\cC^{\infty}$ 类的。
\end{proof}

为了得出矛盾，让我们来证明：满足该引理两个条件的\trVectorField{}必定为零。第二个条件也可以写成
\[
\int_{\RR\times S^{1}}dh_{t}(Z)\,ds\,dt=0\text{。}
\]
当 $Z(s,t)=\lambda(t)\,\partial u/\partial s$ 时，这个关系式是满足的：
\[
\int_{\RR\times S^{1}}dh_{t}\left(\lambda(t)\frac{\partial u}{\partial s}\right)ds\,dt=\int_{S^{1}}\lambda(t)\int_{\RR}\frac{\partial}{\partial s}h_{t}(u(s,t))\,ds\,dt=0
\]
这是因为，按 $\cC^{\infty}_{\varepsilon}(H_{0})$ 的定义，$h$ 在 $X_{H}$ 的 1-\trPeriodic\trTrajectory{}附近为零。反之，我们有下面的结果。

\begin{lemma}\label{thm:z-lambda-grad}
若 $Z\in L^{q}(\RR\times S^{1};u^{\star}TW)$ 是 $\cC^{\infty}$ 类的，并且满足
\[
\forall h\in\cC^{\infty}_{\varepsilon}(H_{0})\text{，}\quad\langle Z,\operatorname{grad}_{u}h\rangle=0\text{，}
\]
那么存在一个 $\cC^{\infty}$ 函数 $\lambda\colon S^{1}\to\RR$，使得
\[
Z(s,t)=\lambda(t)\frac{\partial u}{\partial s}(s,t)\text{。}
\]
\end{lemma}

\begin{proof}
这个证明依赖于\trFloer、\trHofer{}与\trSalamon{}的一条定理，该定理保证\trFloerEquation{}的每一个解都是“\trSomewhereInjective”的——这个概念的具体细节将在接下来的几个定义以及\cref{thm:critical-points-discrete}中给出；该性质本身则将在\cref{sec:somewhere-injective}中证明。

设 $u\in\scr{M}(x,y,J,H)$。称 $(s_{0},t_{0})\in\RR\times S^{1}$ 是 $u$ 的一个\trCriticalPoint，若 $\partial u/\partial s(s_{0},t_{0})=0$。一个点 $(s_{0},t_{0})$ 称为\trRegular{}的，如果它不是\trCriticalPoint，并且对每个 $s\ne s_{0}$ 都有 $u(s_{0},t_{0})\ne u(s,t_{0})$（即使 $s=\pm\infty$ 也不例外，也就是说，$u(s_{0},t_{0})\ne x(t_{0})$ 或 $u(s_{0},t_{0})\ne y(t_{0})$）。\trCriticalPoint{}所成之集记作 $C(u)\subset\RR\times S^{1}$。我们令 $R(u)$ 表示 $u$ 的\trRegular{}点所成之集。

\begin{theorem}[{[32]}]\label{thm:critical-points-discrete}
$C(u)$ 这个\trCriticalPoint{}之集是离散的，而 $R(u)$ 这个\trRegular{}点之集是 $\RR\times S^{1}$ 的一个\trDense\trOpen{}子集。
\end{theorem}

我们将在\cref{sec:somewhere-injective}中证明这条定理。

设 $(s_{0},t_{0})\in R(u)$。假设 $\partial u/\partial s(s_{0},t_{0})$ 与 $Z(s_{0},t_{0})$ 线性无关。我们将要构造一个 $h\in\cC^{\infty}_{\varepsilon}(H_{0})$，使得 $\langle Z,\operatorname{grad}_{u}h\rangle\ne 0$，从而得出矛盾。我们定义一个映射
\[
\begin{aligned}
\widetilde{u}\colon\RR\times S^{1}&\longrightarrow W\times S^{1}\\
(s,t)&\longmapsto(u(s,t),t)\text{。}
\end{aligned}
\]
我们把 $Z$ 看成 $W\times S^{1}$ 上沿 $\widetilde{u}$ 的一个\trVectorField（它在方向 $\partial/\partial t\in TS^{1}$ 上的分量为零）。我们把 $h$ 看成一个定义在 $W\times S^{1}$ 上的映射。我们写
\[
\int_{\RR\times S^{1}}dh_{t}(Z)\,ds\,dt=\int_{\RR\times S^{1}}dh(Z)\,ds\,dt\text{，}
\]
并构造一个 $h\in\cC^{\infty}_{\varepsilon}(H_{0})$，使得这个积分不为零。注意，只需在 $\widetilde{u}$ 的像的邻域内、沿\trVectorField{} $Z$ 的\trIntegralCurve{}定义 $h$ 就够了。

我们选取一个实数 $\delta>0$，它取得充分小，使得正方形
\[
C_{\delta}=\bigl\{(s,t)\in\RR\times S^{1}\ \bigm|\ |s-s_{0}|\le\delta\ \text{并且}\ |t-t_{0}|\le\delta\bigr\}
\]
含于 $u$ 的由\trRegular{}点\footnote{译注：原书此处印作“regular values of $u$”（$u$ 的\trRegularValue）；按 $R(u)$ 前文方定义为 $u$ 的\trRegular{}点之集（其元素为定义域中的点 $(s,t)$ 而非取值），疑应为 regular points（\trRegular{}点），此处径改。——译者注}所成的\trDense\trOpen{}集 $R(u)$ 之中。于是
\[
\widetilde{u}\colon C_{\delta}\longrightarrow W\times S^{1}
\]
是一个\trEmbedding（倘若 $\delta$ 取得充分小，而它正是如此）。我们还可以——并且也确实——（通过减小 $\delta$）假设：当 $(s,t)\in C_{\delta}$ 时，$Z(s,t)$ 与 $\partial u/\partial s(s,t)$ 作为 $T_{u(s,t)}W$ 中的向量是线性无关的。放在 $W\times S^{1}$ 中看，这个条件就转化成：$Z(s,t)$ 与切于像 $\widetilde{u}(C_{\delta})$ 的平面\trTransverse。

由于 $\widetilde{u}\colon C_{\delta}\to W\times S^{1}$ 是\trEmbedding，必要时再把 $\delta$ 减小一些，我们就可以考虑一个局部\trParametrization
\[
\psi\colon C_{\delta}\times D_{\varepsilon}\longrightarrow W\times S^{1}
\]
（这里 $D_{\varepsilon}=\bigl\{y\in\RR^{2n-1}\bigm|\|y\|\le\varepsilon\bigr\}$），使得 $\psi|_{C_{\delta}\times\{0\}}=\widetilde{u}$。我们用 $B$（取“box”（盒子）之意）来表示 $\psi$ 的像。这个盒子展示于\cref{fig:transversality-boxes}的左边。

\begin{figure}[htbp]
\centering
\begin{tikzpicture}[line width=0.6pt,scale=0.92]
  \begin{scope}
    \draw (0,0) rectangle (3.4,2.1);
    \draw (3.4,2.1) -- (4.2,2.7) (3.4,0) -- (4.2,0.6);
    \draw[dashed] (0,2.1) -- (0.8,2.7) (0,0) -- (0.8,0.6);
    \draw (0.8,0.6) -- (4.2,0.6) (0.8,0.6) -- (0.8,2.7) (0.8,2.7) -- (4.2,2.7);
    \draw (0,1.05) -- (3.4,1.05) -- (4.2,1.65) -- (0.8,1.65) -- cycle;
    \node at (4.6,1.35) {$\widetilde{u}(C_{\delta})$};
    \foreach \x/\y in {0.30/0.62,0.62/0.72,0.94/0.60,1.26/0.70,1.58/0.62,1.90/0.72,2.22/0.62,2.54/0.70,2.86/0.60,3.10/0.70} {
      \draw[->] (\x,\y) -- (\x+0.20,\y+0.78);
    }
    \node at (1.15,2.62) {$Z$};
  \end{scope}
  \begin{scope}[xshift=7.6cm]
    \draw (0,0) rectangle (3.4,1.9);
    \draw (0,1.9) -- (1.1,3.04) (3.4,1.9) -- (4.5,3.04) (3.4,0) -- (4.5,1.14);
    \draw (1.1,3.04) -- (4.5,3.04) -- (4.5,1.14);
    \draw (2.8,1.14) -- (4.5,1.14);%
    \draw[dashed] (1.1,3.04) -- (1.1,2.48);%
    \draw (1.1,2.48) -- (2.8,1.14) -- (1.7,0) -- (0,1.35) -- cycle;
    \draw (0,0) -- (0.75,0.75);
    \draw (1.1,2.48) .. controls (0.65,2.9) and (0.1,2.8) .. (-0.25,2.5)
          .. controls (-0.6,2.2) and (-0.72,2.0) .. (-0.82,1.68)
          .. controls (-1.05,1.35) and (-1.3,0.95) .. (-1.15,0.55)
          .. controls (-1.0,0.2) and (-0.5,0.05) .. (0,0);
    \draw (-0.82,1.6) .. controls (-0.55,1.68) and (-0.3,1.45) .. (0,1.35);
    \node at (-0.68,3.12) {$\mathrm{Im}(\widetilde{u})$ 的一部分};
    \draw[->] (2.34,0.67) .. controls (2.37,1.05) .. (2.44,1.79);
    \node at (3.1,1.7) {$Z(s,t)$};
    \draw[->] (2.36,0.68) .. controls (2.42,0.28) .. (2.5,-0.13);
    \node at (3.35,-0.38) {$Z(s',t')$};
  \end{scope}
\end{tikzpicture}
\caption{}
\label{fig:transversality-boxes}
\end{figure}

我们将要构造一个\trSupport{}含于 $B$ 中的 $h$，办法是预先指定它沿 $\widetilde{u}(C_{\delta})$ 在\trVectorField{} $Z$ 的\trIntegralCurve{}上的取值。倘若圆柱形的像 $\mathrm{Im}(\widetilde{u})$ 在 $\widetilde{u}(C_{\delta})$ 之外与 $B$ 相交（如\cref{fig:transversality-boxes}右边所示），这件事就会变得困难。我们借助一条辅助引理来排除这种可能性。

\begin{lemma}\label{thm:image-box-criterion}
对充分小的 $\varepsilon$，$\widetilde{u}(s,t)\in\mathrm{Im}(\psi)$ 当且仅当 $(s,t)\in C_{\delta}$。
\end{lemma}

\begin{proof}
倘若引理所断言的性质不真，那么就存在一个趋于 $0$ 的序列 $\varepsilon_{n}$ 与一列 $(s_{n},t_{n})\in\RR\times S^{1}-C_{\delta}$，使得
\[
\widetilde{u}(s_{n},t_{n})=\psi(s_{n}',t_{n}',y_{n})\qquad\text{其中 }(s_{n}',t_{n}',y_{n})\in C_{\delta}\times D_{\varepsilon_{n}}\text{，}
\]
并且
\[
\lim_{n\to+\infty}(s_{n}',t_{n}',y_{n})=(s_{\star},t_{\star},0)\in C_{\delta}\times\RR^{2n-1}\text{。}
\]
必要时抽取一个子序列之后，序列 $(s_{n},t_{n})$ 趋于 $(-\infty,t_{\#})$、$(+\infty,t_{\#})$ 或 $(s_{\#},t_{\#})\in\RR\times S^{1}$。我们必有 $t_{\#}=t_{\star}$，而且要么 $x(t_{\star})=u(s_{\star},t_{\star})$、$y(t_{\star})=u(s_{\star},t_{\star})$，要么 $u(s_{\#},t_{\#})=u(s_{\star},t_{\star})$。这三种相等关系中的每一种都与 $(s_{\star},t_{\star})$ 是 $u$ 的一个\trRegular{}点这一事实相矛盾，而它必定是\trRegular{}点，因为 $(s_{\star},t_{\star})\in C_{\delta}\subset R(u)$。
\end{proof}

现在我们可以定义 $h$ 了。在 $B$ 之外，我们取 $h\equiv 0$。注意，这样一来 $h$ 在 $X_{H}$ 的 1-\trPeriodic\trTrajectory{}附近为零。为了在 $B$ 的内部定义 $h$，我们选取 $C_{\delta}$ 上一个 $\cC^{\infty}$ 类的正函数 $\beta$，它在某个 $C_{\delta'}$（$\delta'<\delta$）上取值 $1$，而在 $C_{\delta}$ 的边界的一个邻域中取值零。设 $\gamma_{s,t}\colon[-\varepsilon_{0},\varepsilon_{0}]\to B$ 是一条曲线，使得
\[
\gamma_{s,t}(0)=\widetilde{u}(s,t)\qquad\text{以及}\qquad\gamma_{s,t}'(0)=Z(s,t)\text{。}
\]
公式 $h(\gamma_{s,t}(\sigma))=\sigma\beta(s,t)$ 在盒子的一个与 $[-\varepsilon_{0},\varepsilon_{0}]\times C_{\delta}$ \trDiffeomorphic{}的子集上定义了一个函数，这是因为 $Z$ 与像 $\widetilde{u}(C_{\delta})$\trTransverse。而且，在这个 $B$ 的子集上，我们有
\[
(dh)_{\widetilde{u}(s,t)}(Z(s,t))=\beta(s,t)\text{。}
\]
我们把 $h$ 任意地延拓成一个\trSupport{}含于 $B$ 中的函数。利用辅助引理\cref{thm:image-box-criterion}，我们最终得到
\[
\int_{\RR\times S^{1}}(dh)_{\widetilde{u}}(Z)\,ds\,dt=\int_{C_{\delta}}(dh)_{\widetilde{u}}(Z)\,ds\,dt=\int_{C_{\delta}}\beta(s,t)\,ds\,dt>0\text{。}
\]
由于\trSupport{}含于 $B$ 中的那些 $\cC^{\infty}_{\varepsilon}(H_{0})$ 函数，在\trSupport{}含于 $B$ 中的 $\cC^{\infty}$ 函数所成的空间中是\trDense{}的（命题~8.3.4），我们可以用一个属于 $\cC^{\infty}_{\varepsilon}(H_{0})$ 的函数去逼近 $h$，使得这个积分仍然保持为正。

$\partial u/\partial s(s_{0},t_{0})$ 与 $Z(s_{0},t_{0})$ 线性无关这一假设把我们引到了矛盾。于是，存在一个函数 $\lambda\colon R(u)\to\RR$，使得对每个 $(s,t)\in R(u)$，
\[
Z(s,t)=\lambda(s,t)\frac{\partial u}{\partial s}(s,t)\text{。}
\]
为了结束我们这条引理（即\cref{thm:z-lambda-grad}）的证明，现在还必须证明 $\lambda(s,t)$ 不依赖于 $s$。函数 $\lambda$ 这时就能由公式
\[
\lambda(s,t)=\frac{\bigl\langle Z(s,t),\partial u/\partial s(s,t)\bigr\rangle}{\bigl\|\partial u/\partial s(s,t)\bigr\|^{2}}
\]
表出，这就表明它是 $R(u)$ 上的一个 $\cC^{\infty}$ 函数。这个公式还使我们得以把 $\lambda$ 延拓到 $\RR\times S^{1}$ 中 $C(u)$ 的补集上去。由于 $R(u)$ 在 $\RR\times S^{1}$ 中\trDense，关系式 $Z=\lambda\,\partial u/\partial s$ 在 $\RR\times S^{1}-C(u)$ 上依然成立。对每个 $h\in\cC^{\infty}_{\varepsilon}(H_{0})$，我们有
\[
\begin{aligned}
\int_{\RR\times S^{1}}(dh)_{\widetilde{u}}(Z)\,ds\,dt&=\int_{\RR\times S^{1}}\lambda(s,t)(dh)_{\widetilde{u}}\left(\frac{\partial u}{\partial s}\right)ds\,dt\\
&=\int_{\RR\times S^{1}}\lambda(s,t)\frac{\partial}{\partial s}\bigl(h(u(s,t),t)\bigr)ds\,dt\text{。}
\end{aligned}
\]
假设存在一点 $(s_{0},t_{0})\in\RR\times S^{1}-C(u)$，使得 $\partial\lambda/\partial s$ 在其处不为零。这个性质在以 $(s_{0},t_{0})$ 为中心且含于 $\RR\times S^{1}-C(u)$ 的某个正方形 $C_{\delta}$ 上依然成立。与前面一样，我们构造 $W\times S^{1}$ 中的一个盒子 $B$，使得当 $h$ 的\trSupport{}含于 $B$ 中时，
\[
\int_{\RR\times S^{1}}(dh)_{\widetilde{u}}(Z)\,ds\,dt=\int_{C_{\delta}}(dh)_{\widetilde{u}}(Z)\,ds\,dt\text{。}
\]
我们再一次利用上面选定的\trPlateauFunction{} $\beta\colon C_{\delta}\to\RR$ 来定义
\[
h(\widetilde{u}(s,t))=-\beta(s,t)\frac{\partial\lambda}{\partial s}(s,t)\qquad\text{于 }\widetilde{u}(C_{\delta})\text{ 上}
\]
（并把它任意地延拓成一个\trSupport{}含于 $B$ 中的函数）。于是在 $B$ 中使用分部积分，便得到
\[
\begin{aligned}
\int_{\RR\times S^{1}}(dh)_{\widetilde{u}}(Z)\,ds\,dt&=\int_{C_{\delta}}(dh)_{\widetilde{u}}(Z)\,ds\,dt\\
&=\int_{C_{\delta}}\lambda(s,t)\frac{\partial}{\partial s}\bigl(h(u(s,t),t)\bigr)ds\,dt\\
&=-\int_{C_{\delta}}\frac{\partial\lambda}{\partial s}(s,t)\,h(u(s,t),t)\,ds\,dt\\
&=\int_{C_{\delta}}\beta(s,t)\left(\frac{\partial\lambda}{\partial s}(s,t)\right)^{2}ds\,dt\\
&\ge\int_{C_{\delta'}}\left(\frac{\partial\lambda}{\partial s}(s,t)\right)^{2}ds\,dt>0\text{。}
\end{aligned}
\]
再由\trDense{}性，我们可以——并且也确实——像前面那样假设 $h\in\cC^{\infty}_{\varepsilon}(H_{0})$。这样我们就与\cref{thm:z-pairing-vanish}相矛盾了，于是那个假设（存在一对 $(s_{0},t_{0})$ 使得 $\partial\lambda/\partial s$ 不为零）不成立。所以我们有
\[
\forall(s,t)\in\RR\times S^{1}-C(u)\text{，}\qquad\frac{\partial\lambda}{\partial s}(s,t)=0\text{。}
\]
由于 $C(u)$ 是离散的，$\RR\times S^{1}-C(u)$ 是\trConnected{}的（这个断言留给我们的读者当作习题）。因此在 $\RR\times S^{1}-C(u)$ 上，$\lambda(s,t)=\lambda(t)$ 不依赖于 $s$。我们于是可以把这个函数延拓成 $\RR^{2}$ 上的一个 $\cC^{\infty}$ 函数，只需简单地令
\[
\lambda(s_{0},t_{0})=\lambda(s_{0}+\varepsilon,t_{0})\qquad\text{若 }(s_{0},t_{0})\in C(u)\text{。}
\]
关系式
\[
Z(s,t)=\lambda(t)\frac{\partial u}{\partial s}(s,t)
\]
因而在 $\RR\times S^{1}$ 上成立，\cref{thm:z-lambda-grad}得证。
\end{proof}

\begin{proof}[命题~8.1.4~的证明]
让我们来证明：对每个 $t$ 都有 $\lambda(t)\ne 0$。倘若不然，就会存在一个 $t_{0}$，使得对每个 $s\in\RR$ 都有 $Z(s,t_{0})=0$。特别地，所有的导数 $\partial^{k}Z/\partial s^{k}(s,t_{0})$ 都将为零。由于 $Z$ 是 $L^{\star}Z=0$ 的一个解，我们可以用归纳法证明，$Z$ 的一切导数对一切 $s$ 与 $t=t_{0}$ 都将为零。由
\[
L^{\star}Z=-\frac{\partial Z}{\partial s}+J_{0}\frac{\partial Z}{\partial t}+{}^{t}S(s,t)Z
\]
可以推出，$Z$ 将是“\trPerturbation{}的\trCauchyRiemann”型方程的一个解，从而满足一条\textbf{\trContinuationPrinciple}（continuation principle），它与\trHolomorphic{}映射所满足的那条相类似：如果一个映射与其一切导数在 $\RR\times\{t_{0}\}$ 上都为零，那么它处处为零。我们将会解释“\trPerturbation{}的\trCauchyRiemann{}方程”这一概念，并证明这条\trContinuationPrinciple（\cref{thm:continuation-principle}）。

这条\trContinuationPrinciple{}于是将蕴含 $Z$ 恒为零，这就与我们的假设相矛盾。

因此，对每个 $t$ 都有 $\lambda(t)\ne 0$。我们还可以——并且也确实——进一步假设
\[
\lambda(t)>0\qquad\text{对一切 }t\text{。}
\]
这时就有
\[
\int_{0}^{1}\Big\langle\frac{\partial u}{\partial s}(s,t),Z(s,t)\Big\rangle dt=\int_{0}^{1}\lambda(t)\Big|\frac{\partial u}{\partial s}(s,t)\Big|^{2}dt>0\qquad\text{对一切 }s\in\RR\text{；}
\]
因此，这个关于 $s$ 的函数是正的，并且当 $s\to\pm\infty$ 时趋于 $0$（正如 $\partial u/\partial s$ 那样）。倘若这个函数是常数，就会导致矛盾。所以，让我们最终来证明：它其实确实是常数。回忆一下，正如我们已经指出过的（\cref{thm:df-kernel-nontrivial}），$\partial u/\partial s$ 是 $LY=0$ 的一个解。为了简化记号，我们令
\[
Y(s,t)=\frac{\partial u}{\partial s}(s,t)\text{。}
\]
我们有 $LY=0$ 与 $L^{\star}Z=0$，也就是说，
\[
\frac{\partial Y}{\partial s}=-J_{0}\frac{\partial Y}{\partial t}-SY\qquad\text{以及}\qquad\frac{\partial Z}{\partial s}=J_{0}\frac{\partial Z}{\partial t}+{}^{t}SZ\text{。}
\]
于是，我们的函数关于 $s$ 的导数为
\[
\begin{aligned}
\frac{d}{ds}\int_{0}^{1}\langle Y,Z\rangle\,dt&=\int_{0}^{1}\left(\Big\langle\frac{\partial Y}{\partial s},Z\Big\rangle+\Big\langle Y,\frac{\partial Z}{\partial s}\Big\rangle\right)dt\\
&=\int_{0}^{1}\left(\Big\langle-J_{0}\frac{\partial Y}{\partial t},Z\Big\rangle-\langle SY,Z\rangle+\Big\langle Y,J_{0}\frac{\partial Z}{\partial t}\Big\rangle+\langle Y,{}^{t}SZ\rangle\right)dt\\
&=-\int_{0}^{1}\left(\Big\langle J_{0}\frac{\partial Y}{\partial t},Z\Big\rangle+\Big\langle J_{0}Y,\frac{\partial Z}{\partial t}\Big\rangle\right)dt\\
&=-\int_{0}^{1}\frac{\partial}{\partial t}\langle J_{0}Y,Z\rangle\,dt=0\text{。}
\end{aligned}
\]
这个矛盾蕴含 $Z$ 为零，这就结束了命题~8.1.4 证明中的\trSurjective{}部分。
\end{proof}

$\Gamma$ 具有一个连续右逆的证明，由下面的抽象引理给出。

\begin{lemma}\label{thm:continuous-right-inverse}
设 $E$、$F$ 与 $G$ 是\trBanachSpace，而
\[
L_{1}\colon E\to G\text{，}\qquad L_{2}\colon F\to G
\]
是\trLinearOperator。假设 $L_{1}$ 是\trFredholm{}的，并且由 $\Gamma(x,y)=L_{1}(x)+L_{2}(y)$ 定义的 $\Gamma\colon E\oplus F\to G$ 是\trSurjective{}的。那么，$\Gamma$ 具有一个连续右逆。
\end{lemma}

\begin{proof}
写 $G=\mathrm{Im}(L_{1})\oplus H$，其中 $H$ 是闭的且有限维。设 $E'$ 是 $E$ 的一个闭子空间，使得 $E=\mathrm{Ker}(L_{1})\oplus E'$。算子 $L_{1}\colon E'\to\mathrm{Im}(L_{1})$ 显然是\trBijective{}的。把复合
\[
L_{1}^{-1}\colon\mathrm{Im}(L_{1})\longrightarrow E'\subset E\subset E\oplus F
\]
记作 $L_{1}^{-1}$。设 $h_{1},\dots,h_{r}$ 是 $H$ 的一组基，并设 $x_{1},\dots,x_{r}\in E\oplus F$ 满足 $\Gamma(x_{i})=h_{i}$。用 $\nu(h_{i})=x_{i}$ 来定义 $\nu\colon H\to E\oplus F$。这是一个连续映射，因为它的像是有限维的。现在，映射
\[
\begin{aligned}
\Pi\colon\mathrm{Im}(L_{1})\oplus H&\longrightarrow E\oplus F\\
(z,h)&\longmapsto(L_{1}^{-1}(z),0)+\nu(h)
\end{aligned}
\]
是 $\Gamma$ 的一个右逆：容易验证 $\Gamma\circ\Pi=\mathrm{Id}$；此外，$\Pi$ 是连续的，因为它可以写成
\[
\Pi=(L_{1}^{-1}\circ\mathrm{pr}_{\mathrm{Im}(L_{1})},0)+\nu\circ\mathrm{pr}_{H}\text{。}
\]
\end{proof}

\subsection{定理~8.1.1~与~8.1.2~证明的准备，$\scr{Z}(x,y,J)$ 的\trTangentSpace}
\label{subsec:tangent-space-z}

在前一个证明的末尾所得到的那个矛盾，同样也结束了命题~8.1.3 的证明。我们已经证明了 $\scr{Z}(x,y,J)$ 是一个\trBanachManifold。它在点 $(u,h)$ 处的\trTangentSpace{}是 $(d\sigma)_{(u,h)}$ 的\trKernel，也就是说，是空间
\[
\bigl\{(Y,\eta)\in W^{1,p}(\RR\times S^{1};\RR^{2n})\times\cC^{\infty}_{\varepsilon}(H_{0})\ \bigm|\ (d\scr{F})_{u}Y+\operatorname{grad}_{u}\eta=0\bigr\}\text{。}
\]
我们考虑投影
\[
\begin{aligned}
\pi\colon\scr{Z}(x,y,J)&\longrightarrow\cC^{\infty}_{\varepsilon}(H_{0})\\
(u,H_{0}+h)&\longmapsto h
\end{aligned}
\]
及其\trTangentMap
\[
\begin{aligned}
(d\pi)_{(u,H)}\colon T_{(u,H)}\scr{Z}(x,y,J)&\longrightarrow T_{h}\cC^{\infty}_{\varepsilon}(H_{0})=\cC^{\infty}_{\varepsilon}(H_{0})\\
(Y,\eta)&\longmapsto\eta
\end{aligned}
\]
（其中 $H=H_{0}+h$）。注意，$\pi$ 是\trSmooth{}的。让我们来验证 $\pi$ 是一个\trFredholm{}映射。$(d\pi)_{(u,H_{0}+h)}$ 的\trKernel{}恰好就是 $L=(d\scr{F})_{u}$ 的\trKernel，因而维数有限。它的\trImage{}是 $L$ 的\trImage{}在线性映射
\[
\begin{aligned}
\cC^{\infty}_{\varepsilon}(H_{0})&\longrightarrow L^{p}(u^{\star}TW)\\
\eta&\longmapsto\operatorname{grad}_{u}\eta
\end{aligned}
\]
之下的\trPreimage，从而是余维数有限的（并且是闭的）。

现在我们可以应用\trSard--\trSmale{}定理了，这条定理如下。

\begin{theorem}[\trSard--\trSmale~{[73]}]\label{thm:sard-smale}\BookIndexDefinition{T012}
设 $E$ 与 $F$ 是两个\trSeparable{}的\trBanachSpace，设 $U$ 是 $E$ 的一个\trOpen{}子集，并设
\[
\scr{F}\colon U\longrightarrow F
\]
是一个（\trSmooth{}的）\trFredholm{}映射。那么，$\scr{F}$ 的\trRegularValue{}所成之集\footnote{译注：原书此处集合记号印作“$y=f(x)$”；按本定理中的映射记作 $\scr{F}$，疑 $f$ 为 $\scr{F}$ 之误，此处径改。——译者注}
\[
\bigl\{y\in F\ \bigm|\ \text{若 }y=\scr{F}(x)\text{ 对某个 }x\in U\text{ 成立，则 }T_{x}\scr{F}\text{ 是\trSurjective{}的}\bigr\}
\]
是可数个\trDense\trOpen{}子集的交。
\end{theorem}

\begin{remark}\label{thm:sard-smale-remark}
在这条定理中，\trSeparable{}性保证我们可以从任何\trOpen{}\trCovering{}中抽出一个可数子\trCovering。要想得到\emph{可数}个\trDense\trOpen{}子集的交、从而得到一个\trDense{}子集，这个条件是必不可少的。
\end{remark}

现在让我们来证明下面的结果。

\begin{lemma}\label{thm:regular-values-pi}
$\pi$ 的\trRegularValue{}恰好是那些使得对每个 $u\in\scr{M}(H_{0}+h,J)$，$(d\scr{F})_{u}$ 都是\trSurjective{}的 $h\in\cC^{\infty}_{\varepsilon}(H_{0})$。
\end{lemma}

\begin{proof}
设 $h$ 是 $\pi$ 的一个\trRegularValue，并设 $u\in\scr{M}(H_{0}+h,J)$。我们将要证明 $(d\scr{F})_{u}$ 是\trSurjective{}的。倘若不然，那么正如在\cref{thm:z-pairing-vanish}中那样，就会存在一个 $Z\in L^{q}(\RR\times S^{1};\RR^{2n})$，使得
\[
\forall Y\in W^{1,p}(\RR\times S^{1};\RR^{2n})\text{，}\quad\langle(d\scr{F})_{u}(Y),Z\rangle=0\text{。}
\]
由于 $(d\pi)_{(u,H)}$ 是\trSurjective{}的，故对每个 $\eta\in\cC^{\infty}_{\varepsilon}(H_{0})$，都存在一个\trVectorField{} $Y\in W^{1,p}(\RR\times S^{1};\RR^{2n})$，使得
\[
(d\scr{F})_{u}(Y)+\operatorname{grad}_{u}\eta=0\text{。}
\]
利用这两个等式，命题~8.1.4 的证明（见\cref{subsec:proof-prop-814}）就给出 $Z=0$，遂得矛盾。因此，$(d\scr{F})_{u}$ 是\trSurjective{}的。

反之，设 $h\in\cC^{\infty}_{\varepsilon}(H_{0})$，并设对每个 $u\in\scr{M}(H_{0}+h,J)$，$(d\scr{F})_{u}$ 都是\trSurjective{}的。那么，对给定的 $\eta\in\cC^{\infty}_{\varepsilon}(H_{0})$，选取一个 $Y\in W^{1,p}(\RR\times S^{1};\RR^{2n})$，使得
\[
(d\scr{F})_{u}(Y)=-\operatorname{grad}_{u}\eta\text{。}
\]
于是 $(Y,\eta)\in T_{(u,H)}\scr{Z}(x,y,J)$，并且 $(d\pi)_{(u,H)}(Y,\eta)=\eta$，从而 $(d\pi)_{(u,H)}$ 是\trSurjective{}的，而 $h$ 就是 $\pi$ 的一个\trRegularValue。
\end{proof}

\cref{thm:regular-values-pi}立即给出定理~8.1.1（这里用到\trSard--\trSmale{}定理，也就是\cref{thm:sard-smale}）。为了使这个证明真正完备，我们仍须证明\cref{thm:critical-points-discrete}与\trContinuationPrinciple（它将作为\cref{thm:continuation-principle}而得到证明）。

一旦承认了这些结果，让我们最终来证明（必要时在原\trHamiltonian{}上取一个\trPerturbation{}之后）：从 $x$ 到 $y$ 的\trTrajectory{}所成的空间是一个具有所要维数的\trManifold。换句话说，让我们来证明定理~8.1.2。

\begin{proof}[定理~8.1.2~的证明]
设 $h\in\scr{H}_{\mathrm{reg}}$。那么由\cref{thm:regular-values-pi}，$h$ 是
\[
\pi\colon\scr{Z}(x,y,J)\longrightarrow\cC^{\infty}_{\varepsilon}(H_{0})
\]
的一个\trRegularValue。于是，$\pi^{-1}(h)$ 是一个\trManifold，并且它的维数等于 $\pi$ 的\trFredholm{}\trIndex，即
\[
\begin{aligned}
\dim\mathrm{Ker}(d\pi)_{(u,H)}&=\dim\mathrm{Ker}(d\scr{F})_{u}\\
&=\operatorname{Ind}(d\scr{F})_{u}\\
&=\mu(x)-\mu(y)\text{。}
\end{aligned}
\]
$\pi^{-1}(h)$ 的元素是\footnote{译注：原书此处印作“$(H_{0}+h_{0},J)$”；按 $\pi^{-1}(h)$ 的元素应为 $H_{0}+h$ 的\trFloerEquation{}之解（紧随其后的 $\pi^{-1}(h)\subset\scr{M}(x,y,H_{0}+h)$ 亦然），疑 $h_{0}$ 为 $h$ 之误，此处径改。——译者注} $(H_{0}+h,J)$ 的\trFloerEquation{}在 $\scr{P}^{1,p}$ 中的解。由\trEllipticRegularity{}（参见命题~12.1.4），我们有
\[
\pi^{-1}(h)\subset\scr{M}(x,y,H_{0}+h)\text{。}
\]
命题~8.2.3 给出了另一个方向的包含关系，于是
\[
\pi^{-1}(h)=\scr{M}(x,y,H_{0}+h)\text{，}
\]
这正是所要证明的。
\end{proof}

\begin{remark}\label{thm:moduli-topology-remark}
我们已经证明了：对 $h\in\scr{H}_{\mathrm{reg}}$ 与 $\scr{A}_{H_{0}}$ 的一切\trCriticalPoint{} $x$ 与 $y$，空间 $\scr{M}(x,y,H_{0}+h,J)$ 是 $\scr{P}^{1,p}(x,y)$ 的一个\trSubmanifold。在 $\scr{M}(x,y,H_{0}+h,J)$ 上，我们使用 $\cC^{\infty}_{\mathrm{loc}}$ \trTopology。的确，由于 $p>2$，存在从 $\scr{P}^{1,p}(x,y)$ 到连续函数空间的一个\trInjective，它对于 $\cC^{0}_{\mathrm{loc}}$ \trTopology{}是连续的，而且，正如我们已经在\cref{thm:solution-smoothness-regularity}中指出的，在 $\scr{M}(H_{0}+h,J)$（$\scr{M}(x,y,H_{0}+h,J)$ 是它的一个子空间）上，$\cC^{0}_{\mathrm{loc}}$ 与 $\cC^{\infty}_{\mathrm{loc}}$ 这两个\trTopology{}相重合。
\end{remark}

其存在性曾在定理~8.1.1 中预告过的那个\trDense\trOpen{}集 $\scr{H}_{\mathrm{reg}}$，是把考虑各个不同的 1-\trPeriodic\trTrajectory{}对 $(x,y)$ 所得到的那些不同的\trDense\trOpen{}子集取交而得到的\footnote{译注：本句称 $\scr{H}_{\mathrm{reg}}$ 及诸\trTrajectory{}对各自的子集为\trDense\trOpen，系照原书排印；但按定理~8.1.1 与\trSard--\trSmale{}定理（\cref{thm:sard-smale}），每一对 $(x,y)$ 所得的也只是可数个\trDense\trOpen{}子集之交（可数交一般非开），$\scr{H}_{\mathrm{reg}}$ 作为各\trTrajectory{}对所得诸集之交亦是如此，本处并无升级为\trOpen{}集的论证。严格言之，应作“残余（\trDense{} $G_{\delta}$）集／各\trTrajectory{}对所得的残余集”。——译者注}。

为了结束定理~8.1.5 与定理~8.1.1 的证明，我们还须证明 $(d\scr{F})_{u}$ 确实具有\trFredholm{}性质并计算它的\trIndex，这一点我们将在\cref{sec:fredholm-property}与 §8.8 中进行。
